\begin{abstract}
Large-scale unsupervised language models (LMs) acquire extensive world knowledge and some capacity for reasoning, yet exerting fine-grained control over their outputs remains challenging, largely due to the fully unsupervised approach underpinning their training. To address this limitation, current strategies employ human-annotated data indicating the relative quality of model outputs, subsequently adjusting the LM to better reflect these preferences, frequently through reinforcement learning from human feedback (RLHF). Nonetheless, RLHF is a multi-stage and often delicate process: it requires first training a reward model to represent human judgments, and then applying reinforcement learning to optimize the original unsupervised LM for the learned reward, all while ensuring minimal deviation from the base model. This paper proposes a novel reward model parameterization within RLHF that permits analytical recovery of the corresponding optimal policy, thereby recasting the canonical RLHF objective as a tractable classification loss. We introduce \textit{Direct Preference Optimization} (DPO), an algorithm that is robust, efficient, and computationally straightforward—dispensing with the need for language model sampling or extensive hyperparameter search during fine-tuning. Empirical results demonstrate that DPO is capable of adapting LMs to human preferences with efficacy matching or surpassing established approaches. In particular, DPO outperforms PPO-based RLHF for sentiment control tasks and achieves equal or improved results for summarization and single-turn conversation, all while being markedly easier to use and more resource-efficient.
\end{abstract}

\section{Introduction}
Recent advances demonstrate that large-scale unsupervised language models (LMs) trained on massive datasets exhibit remarkable and sometimes unexpected abilities~\citep{chowdhery2022palm, brown2020language, touvron2023llama, bubeck2023sparks}. Nevertheless, the training data for these models is human-generated and reflects a diverse array of intentions, competencies, and priorities. Not all of these are appropriate to emulate; for instance, although it is beneficial for an AI coding assistant to \textit{recognize} frequent programming errors for the purpose of correction, it is preferable for code generation to emphasize the (perhaps infrequent) demonstrations of exceptional programming skill present within its training corpus. In a similar vein, while the model should \textit{know} about widely held misconceptions, such as those believed by 50\% of people, it is undesirable for the model to propagate these inaccuracies at the same rate in response to related queries. Essentially, the task of narrowing the model's \emph{preferred responses and behaviors} from its extensive repertoire of \textit{knowledge and skills} is fundamental for developing AI systems that are reliable, high-performing, and easy to control \citep{ouyang2022training}. Whereas most contemporary approaches employ reinforcement learning (RL) to align LMs with human preferences, we demonstrate that the commonly used RL-driven objective for preference alignment can, in fact, be solved exactly using a straightforward binary cross-entropy formulation, thereby greatly streamlining the preference alignment workflow.

At a broader level, prevalent techniques embed target behaviors within a language model through handpicked datasets encoding human judgments about what constitutes safe and beneficial conduct. This process of learning human preferences typically follows a substantial phase of unsupervised pre-training on expansive text corpora. While supervised fine-tuning on exemplary human-generated responses offers a direct path for preference shaping, methods leveraging reinforcement learning from human or AI feedback (RLHF/RLAIF; \citep{christiano2017deep,bai2022constitutional}) have proven especially effective. RLHF methods involve constructing a reward model based on collected human preference data and subsequently optimizing the LM policy with RL to favor outputs rated highly by the reward model, while restraining deviation from the pretrained policy. Although RLHF achieves significant success in generating high-quality dialogue and code, its application introduces considerable complexity relative to supervised techniques, demanding the training of several language models and extensive on-policy sampling during training, which results in elevated computational requirements.

This work presents a method for optimizing language models to better conform to human preferences, doing so without explicit reward modeling or the use of reinforcement learning. We introduce \textit{{Direct Preference Optimization} (DPO)}, an algorithm that, while implicitly maximizing the same objective as established RLHF methods (i.e., maximizing reward under a KL-divergence constraint), is both easy to implement and straightforward to train. The DPO algorithm intuitively updates the model by raising the relative log-likelihood of preferred responses compared to less-preferred ones. Importantly, it includes an adaptive, example-specific importance weighting mechanism to avoid the mode collapse issues observed when using a naïve probability ratio formulation. In common with prior approaches, DPO adopts a formal preference framework (for instance, the Bradley-Terry model; \cite{bradley1952rankanalysis}) to gauge the concordance between a candidate reward function and observed human preference labels. Whereas traditional pipelines employ the preference model to craft a loss for reward model training, which is then used to update the policy via reward maximization, DPO leverages a change of variables so that the preference loss is defined directly in terms of the policy. Given annotated data consisting of human judgments between pairs of model responses, DPO allows for straightforward policy optimization through a binary cross-entropy objective, thereby inducing the optimal policy relative to an implicit reward function tailored to the preference dataset.

The principal contribution of this paper is Direct Preference Optimization (DPO), an intuitive, reinforcement learning–free procedure for training language models with human preferences. Experimental results demonstrate that DPO matches or surpasses the effectiveness of existing techniques, including RLHF with PPO, across tasks such as sentiment control, summarization, and conversational modeling, utilizing language models with up to 6B parameters.

Recent advances have demonstrated that self-supervised language models of escalating scale are capable of performing certain tasks in a zero-shot manner \citep{radford2019language} or when provided with only a few examples as prompts \citep{gpt3,megatron,chowdhery2022palm}. Nevertheless, these models' abilities on downstream benchmarks and their capacity to adhere to user intent are often substantially enhanced when further trained on curated datasets containing instructions and corresponding human-generated responses \citep{mishra-etal-2022-cross,sanh2022multitask,chung2022scaling,thoppilan2022lamda}. This method, referred to as ‘instruction-tuning,’ has been shown to support large language models (LLMs) in extending their generalization to novel instructions not present in the tuning data, thereby broadly augmenting model utility \citep{chung2022scaling}. Although instruction-tuning has yielded notable advancements, collecting human judgments that express preferences between pairs of responses is typically more tractable than gathering high-quality demonstrations from experts. In response, later research has concentrated on refining LLMs using datasets reflecting human preferences, resulting in measurable gains in tasks such as translation \citep{kreutzer-etal-2018-reliability}, summarization \citep{stiennon2022learning,ziegler2020finetuning}, narrative generation \citep{ziegler2020finetuning}, and following instructions \citep{ouyang2022training,ramamurthy2023is}. Approaches in this area usually involve two stages: first, they train a neural reward model to match human preferences, typically operationalized through a preference model such as the Bradley-Terry model \citep{bradley1952rankanalysis}; then, they further adapt the language model to maximize this reward, often employing reinforcement learning techniques such as REINFORCE \citep{williams1992reinforce}, proximal policy optimization (PPO; \cite{schulman2017proximal}), or related variants \citep{ramamurthy2023is}. A parallel thread of work exploits LLMs that have been tuned for instruction adherence via human feedback to autonomously generate additional synthetic preference data targeting criteria like safety or harmlessness \citep{bai2022constitutional}, relying primarily on minimal human input in the form of textual evaluation guidelines for the model’s own outputs. Collectively, these methods draw upon and integrate advances from two research domains: reinforcement learning for language model training across diverse objectives~\citep{Ranzato2015SequenceLT,paulus2018a,wu2018learning} and generic strategies for preference-based learning from human feedback \citep{christiano2017deep,kupcsik2018learning}. Although leveraging relative preferences from humans is attractive, the practical difficulties associated with reinforcement learning-based fine-tuning of large language models remain significant; the present study proposes a theoretically grounded method for optimizing based on relative preferences without reliance on reinforcement learning.

Beyond language-focused applications, the task of learning policies based on preference information has been explored within both bandit and reinforcement learning domains, resulting in a variety of methodological developments. One particular line of work, known as contextual dueling bandits (CDB; \cite{yue2012karmed,dudik2015contextual}), deals with contextual bandit problems where action preferences or rankings replace explicit reward signals. In scenarios where concrete rewards are unavailable, CDBs leverage the concept of a \textit{von Neumann winner}: a policy whose expected success rate when compared against any alternative policy meets or exceeds 50\% \citep{dudik2015contextual}. Notably, the CDB paradigm assumes preference feedback is delivered incrementally in an online manner, while human preference learning typically relies on a predetermined offline dataset containing annotated action pairs \citep{yan2022human}. Related to this, \textit{preference-based RL} (PbRL) operates using binary preferences that originate from an underlying, unknown `scoring' function in lieu of standard reward feedback \citep{BusaFekete2014,ruiz2023dueling}. Although there exist several PbRL methods—including those designed for utilizing off-policy preference samples—these strategies usually proceed by first inferring the hidden scoring function (i.e., constructing a reward model) and then conducting optimization over it \citep{jain2013learning,BusaFekete2014,christiano2017deep,sadigh2017active,kupcsik2018learning}. By contrast, our contribution is to introduce a unified approach that directly trains a policy to align with the observed preferences in a single stage.

\section{Preliminaries}\label{section:prelims}

Here, we summarize the RLHF framework as detailed in \citeauthor{ziegler2020finetuning} and subsequently in works such as \citep{stiennon2022learning, bai2022training, ouyang2022training}. This process generally unfolds in three sequential steps: (1) supervised fine-tuning (SFT); (2) collecting preference feedback and building a reward model; and (3) reinforcement learning (RL) policy optimization.

\textbf{SFT}: RLHF commonly starts with supervised training of a pre-existing language model using curated, high-quality data related to the target tasks (for example, dialogue or summarization), yielding an initial model denoted by $\pi^\text{SFT}$.

\textbf{Reward Modelling Phase}: The subsequent stage involves generating response pairs by prompting the SFT model with an input $x$, resulting in pairs $(y_1, y_2)\sim \pi^\text{SFT}(y \mid x)$. These candidate outputs are then assessed by human annotators, who indicate a preferred response: $y_w\succ y_l \mid x$, where $y_w$ and $y_l$ are the more and less favored outputs among $(y_1, y_2)$, respectively. Although the preferences originate from an underlying, latent reward function $r^*(y, x)$ (which is unobservable), they serve as the basis for modeling. Among various modeling options, the Bradley-Terry (BT) framework \cite{bradley1952rankanalysis} is frequently employed; more general ranking approaches like the Plackett-Luce model \citep{plackett1975analysis, luce2012individual} can also be applied when dealing with multiple ranked outputs. Within the BT model, the observed human preference distribution $p^*$ is formalized as follows:

\begin{equation}\label{eq:bradley-terry}
    p^*(y_1\succ y_2 \mid x)=\frac{\exp\left(r^*(x, y_1)\right)}{\exp\left(r^*(x, y_1)\right) + \exp\left(r^*(x, y_2)\right)}.
\end{equation}

Given a static collection of preference comparisons $\mathcal{D}=\bigl\{x^{(i)}, y_w^{(i)}, y_l^{(i)}\bigr\}_{i=1}^N$ sampled from the underlying preference model $p^*$, one can introduce a parameterized reward function $r_{\phi}(x, y)$ and adjust its parameters via maximum likelihood estimation. Casting this as a binary classification task, the associated objective becomes the negative log-likelihood loss:

\begin{equation}\label{eq:reward_model}
    \mathcal{L}_R(r_{\phi}, \mathcal{D}) = -\mathbb{E}_{(x, y_w, y_l)\sim \mathcal{D}}\bigl[\log \sigma(r_{\phi}(x, y_w)- r_{\phi}(x, y_l))\bigr]
\end{equation}

where $\sigma$ denotes the logistic sigmoid. For language model scenarios, the reward network $r_{\phi}(x, y)$ is typically built by augmenting the SFT model $\pi^\text{SFT}(y \mid x)$ with an additional linear transformation applied to the top-level transformer output, yielding a scalar reward prediction~\cite{ziegler2020finetuning}. To encourage stable and well-behaved learning dynamics, it is common to normalize rewards so that  $\mathbb{E}_{x,y\sim \mathcal{D}}\left[r_\phi(x, y)\right] = 0$ across all $x$.

\textbf{RL Fine-Tuning Phase}: In the RL fine-tuning stage, the previously obtained reward model is utilized to guide language model updates. Consistent with the frameworks introduced in~\citep{jaques2017sequence, jaques2020human}, the corresponding optimization task is

\begin{equation}\label{eq:RL}
\max_{\pi_{\theta}}  \mathbb{E}_{x\sim \mathcal{D}, y\sim \pi_{\theta}(y \mid x)}\bigl[r_{\phi}(x, y)\bigr] - \beta\mathbb{D}_{\textrm{KL}}\bigl[\pi_{\theta}(y\mid x)\mid \mid \pi_\text{ref}(y\mid x)\bigr],
\end{equation}

where $\beta$ serves as a regularization coefficient moderating divergence from a reference policy $\pi_\text{ref}$, generally chosen to be the initial SFT model $\pi^\text{SFT}$. The trainable policy $\pi_\theta$ is likewise initialized from $\pi^\text{SFT}$. This KL-based constraint is critical to keep policy updates within the regime where the reward model is trustworthy, while also maintaining answer diversity and avoiding degenerate convergence to a single response mode. Owing to the inherent discreteness in text generation, the objective cannot be differentiated directly and must be addressed through reinforcement learning methods. In standard practice~\citep{ziegler2020finetuning, stiennon2022learning, bai2022training, ouyang2022training}, the objective is written as ${r(x, y) = r_{\phi}(x, y) -\beta (\log \pi_{\theta}(y\mid x) - \log \pi_\text{ref}(y\mid x))}$, which is then optimized using PPO~\cite{schulman2017proximal}. 

\section{Direct Preference Optimization}\label{sec:DPO}

In view of the complexities associated with large-scale reinforcement learning for language model adaptation, our objective is to formulate a more straightforward policy optimization technique that acts directly on preference data. Unlike conventional RLHF methods—which rely on first estimating a reward model, followed by RL-driven policy refinement—our method introduces a special reward parameterization enabling direct, closed-form derivation of the optimal policy, thereby circumventing the need for an RL training phase. In detail, our central contribution is to exploit an analytical relationship linking reward models to their corresponding optimal policies, so that a reward-based loss can be re-expressed as a loss over policy parameters. Through this change of variables, we obviate the need for training a separate reward model, yet remain fully compatible with established human preference models like the Bradley-Terry model. As a result,

\textbf{Derivation of the DPO objective.} We begin from the standard RL objective as in previous literature, Eq.~\ref{eq:RL}, involving a generic reward function $r$. Prior work~\citep{peters2007reinforcement, peng2019advantage, korbak2022reinforcement, go2023aligning} establishes that the solution to the KL-regularized reward maximization problem in Eq.~\ref{eq:RL} is given by:

\begin{equation}\label{eq:op_policy}
    \pi_r(y\mid x) = \frac{1}{Z(x)}\pi_\text{ref}(y\mid x)\exp\left(\frac{1}{\beta}r(x, y)\right),
\end{equation}

with $Z(x) =\sum_{y}\pi_\text{ref}(y\mid x)\exp\left(\frac{1}{\beta}r(x, y)\right)$ as the partition function. Further details can be found in Appendix \ref{app:derivation1}. Estimating the partition function $Z(x)$ is computationally intensive even when the reward is approximated via MLE using $r_{\phi}$ for the true reward $r^*$, as previously discussed~\citep{korbak2022reinforcement, go2023aligning}. Therefore, this particular parameterization is challenging for practical computation. To address this, we can manipulate Eq.~\ref{eq:op_policy} to relate the reward directly to the corresponding optimal policy $\pi_r$, reference policy $\pi_\text{ref}$, and the (unknown) partition function $Z(\cdot)$. Specifically, by applying the logarithm and simplifying, we derive:

\begin{equation}\label{eq:main_eq}
    r(x,y) =\beta \log \frac{\pi_r(y\mid x)}{\pi_\text{ref}(y\mid x)} + \beta \log Z(x).
\end{equation}

This representation can be used for the true reward $r^*$ and the associated optimal policy $\pi^*$. Crucially, the Bradley-Terry model is a function of reward differences between two completions: ${p^*(y_1 \succ y_2 \mid x) = \sigma(r^*(x, y_1) - r^*(x, y_2))}$. Plugging the reparameterization from Eq.~\ref{eq:main_eq} into the preference model in Eq.~\ref{eq:bradley-terry}, we find that the partition function cancels out. This leads to an expression of the preference probability solely in terms of $\pi^*$ and $\pi_\text{ref}$, with no dependence on $Z(x)$. Hence, under the Bradley-Terry model, the optimal RLHF policy $\pi^*$ fulfills:

\begin{equation}\label{eq:objective}
    p^*(y_1\succ y_2 \mid x)=\frac{1}{1 + \exp\left(\beta \log \frac{\pi^*(y_2\mid x)}{\pi_\text{ref}(y_2\mid x)} - \beta \log \frac{\pi^*(y_1\mid x)}{\pi_\text{ref}(y_1\mid x)}\right)}
\end{equation}

See Appendix~\ref{app:derivation2} for a step-by-step derivation. Although Eq.~\ref{eq:objective} is specific to the Bradley-Terry model, an analogous process can be employed to derive corresponding forms for the Plackett-Luce models~\citep{plackett1975analysis, luce2012individual}, which we include in Appendix~\ref{app:plackett_luce_models}.

Given this reformulation, the likelihood of observing human preference data depends directly on the optimal policy rather than the underlying reward function. This enables us to define a maximum likelihood estimation framework for a parameterized policy $\pi_\theta$. By analogy with the approach taken for reward modeling (as shown in Eq.~\ref{eq:reward_model}), the policy optimization problem takes the following form:

\begin{equation}\label{eq:optimum_model}
    \mathcal{L}_\text{DPO}(\pi_{\theta}; \pi_\text{ref}) = -\mathbb{E}_{(x, y_w, y_l)\sim \mathcal{D}}\left[\log \sigma \left(\beta \log \frac{\pi_{\theta}(y_w\mid x)}{\pi_\text{ref}(y_w\mid x)} - \beta \log \frac

In this framework, we model an implicit reward through a different parameterization, for which the associated optimal policy coincides with $\pi_\theta$. Notably, the approach aligns with learning a reparametrized Bradley-Terry model, thereby inheriting favorable theoretical guarantees such as consistency, provided certain assumptions about the distribution of preference data hold \cite{bong2022generalized}. We expand on the theoretical characteristics of DPO and its connections to related literature in Section~\ref{sec:theory}.

\textbf{Mechanistic role of the DPO update.} To elucidate the behavior of DPO, it is instructive to examine the gradient of its loss, $\mathcal{L}_\text{DPO}$. The derivative with respect to the parameters $\theta$ takes the form:

\begin{multline*}\label{eq:gradient}
    \nabla_\theta \mathcal{L}_\text{DPO}(\pi_\theta;\pi_\text{ref}) = \\ -\beta\mathbb{E}_{(x, y_w, y_l) \sim \mathcal{D}} \bigg[\underbrace{\sigma(\hat{r}_\theta(x, y_l) - \hat{r}_\theta (x, y_w))}_\text{places greater emphasis when reward estimates are inaccurate}\bigg[\underbrace{\nabla_\theta\log \pi(y_w \mid x)}_\text{promote $y_w$} - \underbrace{\nabla_\theta\log\pi(y_l \mid x)}_\text{demote $y_l$}\bigg]\bigg],
\end{multline*}

where $\hat{r}_\theta(x, y) = \beta \log \frac{\pi_\theta(y \mid x)}{\pi_\text{ref}(y \mid x)}$ serves as the implicitly derived reward utilizing the current model $\pi_\theta$ alongside a reference model $\pi_\text{ref}$ (refer to Section~\ref{sec:theory} for additional detail). Conceptually, the gradient of $\mathcal{L}_\text{DPO}$ adjusts parameters to raise the probability assigned to preferred continuations $y_w$ while lowering it for less preferred ones $y_l$. Crucially, the magnitude of the update for each example is influenced by the extent to which the implicit reward function $\hat{r}_\theta$ erroneously scores the less preferred completion above the preferred, moderated by $\beta$; in effect, it corrects for misorderings with respect to the KL constraint's strength. Empirically, this reweighting proves critical, as omitting the weighting factor (i.e., applying a simpler variant of this method) often leads to model degradation, as evidenced by results in Appendix Table~\ref{tab:unlikelihood_generations}.

\textbf{DPO overview.} The typical workflow for DPO involves two main stages: First, for each prompt $x$, generate candidate completions $y_1, y_2 \sim \pi_\text{ref}(\cdot \mid x)$ and annotate these with human preference information, forming an offline preference dataset $\mathcal{D} = \{x^{(i)}, y_w^{(i)}, y_l^{(i)}\}_{i=1}^N$. Second, adjust the parameters of the language model $\pi_\theta$ by minimizing the objective $\mathcal{L}_\text{DPO}$, using the provided $\pi_\text{ref}$, $\mathcal{D}$, and target $\beta$ value.

In realistic applications, rather than collecting new samples and preferences, it is desirable to leverage publicly available preference datasets. Since these datasets have typically been constructed using samples from $\pi^\text{SFT}$, we set $\pi_\text{ref} = \pi^\text{SFT}$ if such a model exists. In situations where $\pi^\text{SFT}$ cannot be accessed, we estimate $\pi_\text{ref}$ by fitting it to maximize the likelihood of the preferred responses, i.e., ${\pi_\text{ref} = \argmax_{\pi}\mathbb{E}_{x, y_w \sim \mathcal{D}}\left[\log \pi(y_w \mid x)\right]}$. This approach serves to lessen the effects of any mismatch between the unknown true reference distribution and the employed $\pi_\text{ref}$ within DPO. Additional implementation specifics and details regarding hyperparameters are available in Appendix~\ref{app:implementation}.

\section{Theoretical Analysis of DPO} In the following, we elaborate on DPO's conceptual foundation, deliver theoretical guarantees, and discuss how DPO circumvents difficulties inherent in actor-critic approaches for RLHF (notably, methods like PPO~\cite{schulman2017proximal}).

\label{sec:theory}

\subsection{Your Language Model Is Secretly a Reward Model} DPO accomplishes learning without requiring an explicit reward function or policy optimization through RL, relying instead on a unified maximum likelihood formulation. Observe that the objective defined in Eq.~\ref{eq:main_eq} corresponds to a Bradley-Terry model with the reward defined as $r^*(x, y) = \beta \log\frac{\pi^*_\theta(y \mid x)}{\pi_\text{ref}(y \mid x)}$, and we optimize $\pi_\theta$ accordingly. This mirrors the process of optimizing the reward model as in Eq.~\ref{eq:reward_model}, via a change of variables. Here, we develop the formal underpinnings for this reparameterization, demonstrate its lack of restrictiveness on the expressivity of reward models learned, and show that it enables recovery of the optimal policy. We start by formalizing an equivalence relation over reward functions.

\begin{definition}
We define two reward functions $r(x, y)$ and $r'(x, y)$ as equivalent if and only if ${r(x, y)-r'(x, y) = f(x)}$ holds for some function $f$.     
\end{definition}

It is straightforward to verify that this condition satisfies the properties of an equivalence relation, which thereby partitions the space of reward functions into distinct equivalence classes. The following lemmas summarize key implications of this structure:

\begin{lemma}\label{lemma:same_prefrence} In the context of the Plackett-Luce model, including the Bradley-Terry setting, reward functions belonging to the same equivalence class induce identical preference distributions.
\end{lemma}

\begin{lemma}\label{lemma:same_policy}
    Any pair of reward functions from the same equivalence class result in the same optimal policy under a constrained reinforcement learning formulation.
\end{lemma}

Given the directness of these results, we relegate their proofs to Appendix \ref{app:lemma1}. The first lemma reflects the well-documented phenomenon of under-identification inherent to models within the Plackett-Luce family \cite{plackett1975analysis}. As a consequence of this lack of uniqueness, additional identifiability criteria are commonly imposed to facilitate guarantees regarding the maximum likelihood estimation as in Eq. \ref{eq:reward_model} \cite{bong2022generalized}. The second lemma articulates that any reward function selected from an equivalence class yields an identical optimal policy. Thus, for the purposes of policy optimization, it suffices to identify any representative from the appropriate optimal equivalence class. Our main result, proven in Appendix~\ref{app:thm1}, is stated below:

\begin{theorem}\label{thm:main}
    Provided mild conditions, every reward class consistent with the Plackett-Luce framework (including the Bradley-Terry model) admits a representation of the form ${r(x, y) = \beta \log \frac{\pi(y\mid x)}{\pi_\text{ref}(y\mid x)}}$ for some policy model $\pi(y\mid x)$ and fixed reference model $\pi_\text{ref}(y \mid x)$.
\end{theorem}

\begin{sproof}
    Let $r(x, y)$ be any reward function that defines an optimal policy $\pi_r(y \mid x)$ as per Eq. \ref{eq:op_policy}. We aim to demonstrate that an equivalent reward function (i.e., one within the same class as $r$) can be constructed via the stated reparameterization. Define the projection $f$ as  
\begin{equation}
    f(r; \pi_\text{ref}, \beta)(x, y) = r(x, y) - \beta\log\sum_{y}\pi_\text{ref}(y\mid x)\exp\left(\frac{1}{\beta}r(x, y)\right)
\end{equation}
This operator $f$ effectively adjusts the reward by subtracting the (logarithmic) normalization over $y$ given $x$ and $\pi_\text{ref}$, which depends solely on $x$. Thus, $f(r; \pi_\text{ref}, \beta)(x, y)$ remains in the equivalence class of $r(x, y)$. Substituting $r$ with the right-hand side of Eq.~\ref{eq:main_eq} (valid for any reward), we obtain $f(r; \pi_\text{ref}, \beta)(x, y) = \beta \log \frac{\pi_r(y\mid x)}{\pi_\text{ref}(y\mid x)}$. Consequently, this projection produces a canonical representative of the equivalence class in the specified parametric form, showing that our reparameterization maintains generality.
\end{sproof}

Theorem~\ref{thm:main} can be reinterpreted as indicating which specific reward function is chosen by the DPO reparameterization within each equivalence class, namely, the reward function that fulfills:

\begin{equation}\label{eq:lag_p}
     \sum_{y}\underbrace{\pi_\text{ref}(y\mid x)\exp\left(\frac{1}{\beta}r(x, y)\right)}_{=\pi(y\mid x)\text{, as per Thm.~\ref{thm:main} reparam.}} = 1,
\end{equation}

which ensures that $\pi(y\mid x)$ forms a proper probability distribution (all probabilities are non-negative and sum to unity). Furthermore, using Eq.~\ref{eq:op_policy}, Eq.~\ref{eq:lag_p} can be interpreted as representing the partition function associated with the optimal policy determined by the reward $r(x, y)$. The central principle underlying the DPO algorithm is the ability to introduce constraints on the inherently under-determined Plackett-Luce family of preference models—including Bradley-Terry as a special case—so that the class of expressible reward models is preserved, while at the same time making the optimal policy in Eq. \ref{eq:op_policy} analytically manageable for arbitrary prompts $x$.

\subsection{Instability of Actor-Critic Algorithms} The developed framework also sheds light on why conventional actor-critic methods used in RLHF, such as PPO, may exhibit instability. To analyze this, we adhere to the standard RLHF procedure and examine the RL fine-tuning phase described in Section \ref{section:prelims}. This analysis allows us to connect with the control as inference paradigm \cite{levine2018reinforcement} for the constrained reinforcement learning problem introduced in \ref{eq:RL}. Considering a parameterized policy $\pi_{\theta}(y\mid x)$, the typical objective involves minimizing $\mathbb{D}_{\text{KL}}[\pi_{\theta}(y|x) \mid \mid \pi^*(y\mid x)]$, where $\pi^*$ is the optimal policy generated according to Eq. \ref{eq:optimum_model} for the reward $r_{\phi}(y, x)$. Through straightforward manipulation, this optimization is equivalently expressed as:

\begin{equation}\label{eq:AC}
    \max_{\pi_{\theta}}\mathbb{E}_{\pi_{\theta}(y\mid x)}\bigg[\underbrace{r_{\phi}(x, y) -\beta\log\sum_{y}\pi_\text{ref}(y\mid x)\exp\left(\frac{1}{\beta}r_{\phi}(x, y)\right)}_{f(r_{\phi}, \pi_\text{ref}, \beta)} - \underbrace{\beta\log\frac{\pi_{\theta}(y\mid x)}{\pi_\text{ref}(y\mid x)}}_{\text{KL}}\bigg]
\end{equation}

The objective considered here is identical to that targeted in earlier studies 
\citep{ziegler2020finetuning, stiennon2022learning, bai2022training, ouyang2022training}, where the DPO-equivalent reward formulation is applied to the reward function class $r_{\phi}$. Within this framework, the normalization component within $f(r_{\phi}, \pi_\text{ref}, \beta)$ may be viewed as representing the soft value function for the reference policy $\pi_\text{ref}$. Though this normalization does not alter the location of the optimal solution, omitting it can lead to increased variance in the policy gradient, potentially destabilizing the training process. While incorporating a learned value function could address the normalization, this approach is itself challenging to optimize. Prior approaches often circumvented this by calibrating rewards using a human completion baseline—essentially, a single-sample Monte Carlo approximation of the normalization. In contrast, the reparameterization in DPO constructs a reward function that is baseline-independent. 

\section{Experiments}
Here, we provide empirical assessments of DPO's capacity to train policies directly based on preference data. To begin, within a controlled text generation context, we investigate the efficiency with which DPO balances reward maximization and KL-divergence minimization from the reference policy, comparing its behavior with other preference-based algorithms like PPO. We then test DPO on larger-scale models and more complex RLHF scenarios, including tasks involving summarization and dialogue. Our findings suggest that, with minimal hyperparameter tuning, DPO often matches or surpasses robust baselines such as RLHF utilizing PPO, as well as approaches that select the highest-rewarded trajectory out of $N$ samples based on a learned reward. Prior to detailing the empirical outcomes, we outline the experimental design; further specifics can be found in Appendix~\ref{app:exp_details}.

\textbf{Tasks.} We assess algorithm performance across three distinct open-domain text generation challenges. Across all tasks, policies are trained using a preference dataset denoted by $\mathcal{D}=\bigl\{x^{(i)}, y_w^{(i)}, y_l^{(i)}\bigr\}_{i=1}^N$. For \textbf{controlled sentiment generation}, $x$ represents the beginning of a movie review sampled from the IMDb dataset \cite{maas-EtAl:2011:ACL-HLT2011}, and the generated output $y$ is required to convey positive sentiment. To systematically measure performance in this setting, we synthetically produce preference pairs by leveraging a pre-trained sentiment classifier such that $p(\text{positive}\mid x,y_w)>p(\text{positive}\mid x,y_l)$. In the SFT condition, GPT-2-large is further trained until convergence on the training portion of the IMDB review data (additional methodological specifics in App~\ref{app:sentiment_details}). For the \textbf{summarization} task, $x$ consists of a Reddit forum post, and the policy’s objective is to output a concise summary $y$ highlighting the post's main ideas. Consistent with earlier research, we utilize the Reddit TL;DR summarization corpus \citep{volske-etal-2017-tl} paired with human preference data originally collected by \citeauthor{stiennon2022learning}. Our SFT baseline is based on a model that has been fine-tuned with human-authored forum summaries\footnote{\url{https://huggingface.co/CarperAI/openai_summarize_tldr_sft}}, and RLHF is conducted using the TRLX framework \citep{leandro_von_werra_2023_7790115}. The preference data set stems from \citeauthor{stiennon2022learning} but was produced from responses sampled from an alternative SFT model with a similar training regime. Lastly, in the \textbf{single-turn dialogue} scenario, $x$ denotes a user’s query, which can vary from technical questions about astrophysics to requests for advice. Here, the policy should generate a response $y$ that is both informative and engaging. We draw on the Anthropic Helpful and Harmless dialogue corpus \citep{bai2022training}, comprising 170k dialogues between humans and an AI assistant, where each dialogue concludes with two model-generated responses and an associated human preference label. In this case, no dedicated SFT model exists, so we create an SFT baseline by fine-tuning a general-purpose language model solely on the responses preferred by human annotators.

\textbf{Evaluation.} We utilize two primary evaluation methodologies in our experiments. In a controlled sentiment generation scenario, where the true reward function (implemented as a sentiment classifier) is accessible, we assess each algorithm by examining its trade-off curve between the obtained reward and the KL-divergence relative to the reference policy—an analysis feasible only with ground-truth reward information. In contrast, since direct access to the true reward is unavailable in real-world conditions, we instead measure each algorithm’s \textit{win rate} compared to a baseline policy. Here, GPT-4 is employed as a proxy for human evaluators to judge summary quality and response helpfulness for summarization and single-turn dialogue tasks, respectively. For the summarization task, baseline comparisons use reference summaries from the test set; for dialogue, the baseline consists of the preferred response from the test dataset. Though prior work suggests that large language models can outperform traditional metrics as automated evaluators \citep{Chen2023ExploringTU}, we further substantiate our use of GPT-4 through a dedicated human study (Sec.~\ref{sec:human-judgments}). Our findings show that GPT-4's assessments display a high degree of concordance with human judgments, often matching or surpassing the level of agreement found between human annotators.

\textbf{Methods.} Beyond DPO, we benchmark various established techniques for aligning language models with human preferences. For the summarization task, we use zero-shot prompting with \textbf{GPT-J} \citep{gpt-j}, whereas for single-turn dialogue, we utilize 2-shot prompting with \textbf{Pythia-2.8B} \citep{biderman2023pythia}. We also include the \textbf{SFT} model and introduce \textbf{Preferred-FT}, a supervised fine-tuned model trained to output the preferred completion $y_w$—derived either from the SFT model (for sentiment control and summarization) or a generic language model (for dialogue). Additionally, we incorporate the pseudo-supervised \textbf{Unlikelihood} approach~\citep{welleck2019neural}, which trains the policy to increase the likelihood of $y_w$ while simultaneously \textit{decreasing} the probability assigned to $y_l$; this is modulated via an optional coefficient $\alpha\in[0,1]$ applied to the unlikelihood component. Our analysis also encompasses \textbf{PPO} \citep{schulman2017proximal} utilizing a reward model trained on preference data, as well as \textbf{PPO-GT}, which acts as an oracle by learning from the ground-truth reward in the sentiment-controlled scenario. For sentiment experiments, PPO-GT is implemented both through a standard reference version \cite{leandro_von_werra_2023_7790115} and a customized variant with reward normalization and fine-tuned hyperparameters—optimizations also applied to the standard PPO when used with learned rewards. Lastly, we consider the \textbf{Best of $N$} approach: this baseline samples $N$ candidate responses from either the SFT model (or Preferred-FT for dialogue), returning the response with the highest score per a reward model trained from the preference dataset. While this technique achieves strong empirical results, its reliance on sampling $N$ completions for each input at evaluation renders it computationally costly, limiting its practicality as $N$ increases.

\subsection{How well can DPO optimize the RLHF objective?}

The typical objective in RLHF algorithms—maximizing reward subject to a KL penalty—aims to strike a balance between improving the expected reward and maintaining proximity to the reference policy. Thus, fair algorithmic comparisons necessitate considering both the attained reward and the degree of divergence as measured by KL; a marginal reward improvement does not justify a substantial increase in KL. Figure~\ref{fig:frontier-tldr-main} depicts the tradeoff curve between reward and KL for different algorithms in a sentiment analysis scenario. For each algorithm, we conduct several training runs, varying policy conservativeness via hyperparameters (for PPO, target KL $\in\{3,6,9,12\}$; for unlikelihood, $\beta \in \{0.05,0.1,1,5\}$ and $\alpha\in\{0.05,0.1,0.5,1\}$; for preferred-FT, we use different random seeds), culminating in 22 total experimental runs. Evaluation occurs after every 100 training steps and continues until convergence; we assess each trained policy on a set of held-out prompts, recording both the mean reward from the true reward function and the average sequence-level KL\footnote{Specifically, the sum of per-timestep KL-divergence values.} relative to the reference policy $\text{KL}\left(\pi\mid \mid \pi_\text{ref}\right)$. The results indicate that DPO establishes a substantially superior frontier, attaining the highest levels of reward without incurring excessive KL divergence. This finding is especially striking for two reasons. First, although DPO and PPO share the same underlying optimization target, DPO attains greater efficiency; the reward-KL performance of DPO consistently exceeds that of PPO. Second, DPO surpasses PPO in its frontier performance, \emph{even when PPO is provided with access to ground truth rewards} (PPO-GT).

\subsection{Can DPO scale to real preference datasets?}
\label{sec:dpo-real-datasets}
We proceed to assess the fine-tuning capabilities of DPO on tasks involving summarization and single-turn dialogue. In the context of summarization, commonly used automated metrics like ROUGE often show weak alignment with human preference judgments~\citep{stiennon2022learning}, and previous studies have demonstrated that optimizing LMs with PPO on human-labeled preferences leads to the production of more preferred summaries. Our comparison of methods entails sampling outputs from the test partition of the TL;DR summarization dataset and calculating their mean win rate when matched against reference outputs from the test set. Output samples from every method are drawn using temperature values ranging from 0.0 to 1.0, with corresponding win rates visualized in Figure~\ref{fig:frontier-tldr-main} (right). DPO, PPO, and Preferred-FT all use the identical GPT-J SFT model for fine-tuning\footnote{\url{https://huggingface.co/CarperAI/openai_summarize_tldr_sft}}. Our results indicate that at temperature 0.0, DPO achieves a win rate near 61\%, outperforming PPO, which attains approximately 57\% under optimal temperature conditions. Furthermore, DPO attains a greater peak win rate than the best of $N$ strategy. It is important to highlight that we did not extensively optimize DPO’s $\beta$ hyperparameter, suggesting that these outcomes may not represent its full capacity. Additionally, DPO exhibits considerably more stable performance across different sampling temperatures in contrast to PPO, which can regress to baseline GPT-J behavior at elevated temperatures. Preferred-FT, on the other hand, offers little improvement beyond the SFT baseline. For a direct comparison, Section~\ref{sec:human-judgments} details human assessment, where samples from DPO at temperature 0.25 were favored 58\% of the time over PPO samples at the same temperature.

For single-turn dialogue evaluation, we assess various approaches on a portion of the test split from the Anthropic HH dataset \citep{bai2022training}, focusing exclusively on interactions that involve a single exchange between human and assistant. In the case of GPT-4, we calculate win rates for each method by comparing outputs to the test split's preferred completions. Since a standardized SFT model is lacking for this context, our process begins with a Pythia-2.8B pre-trained model. We fine-tune a reference model using Preferred-FT on the selected completions to ensure they remain consistent with the model's output distribution, and subsequently train with DPO. Our analysis also includes the strongest result from 128 Preferred-FT completions (as performance on the Best of $N$ baseline saturates at $N=128$; see Appendix Figure~\ref{fig:best-of-n}) as well as a two-shot prompt version of the original Pythia-2.8B model. In both settings, DPO matches or surpasses other techniques at optimal temperature values. Additionally, we examine an RLHF model trained via PPO on the Anthropic HH dataset \footnote{\url{https://huggingface.co/reciprocate/ppo_hh_pythia-6B}} from an established repository \footnote{\url{https://github.com/CarperAI/trlx/tree/main/examples/hh}}, but are unable to identify prompts or temperature settings that yield performance superior to the base Pythia-2.8B model. Drawing on our observations from TL;DR and noting that both PPO and Best of 128 target identical reward objectives, we use the Best of 128 approach as an approximate stand-in for PPO-level outcomes. In summary, DPO stands out as the only method that is both computationally efficient and capable of outperforming the preferred completions within the Anthropic HH dataset, achieving results comparable to or exceeding the computationally costly Best of 128 baseline. Notably, Figure~\ref{fig:dialogue-main} indicates that DPO achieves peak performance with relatively rapid convergence.

\subsection{Generalization to a new input distribution}

To more thoroughly investigate how PPO and DPO perform under changes in input distribution, we assessed the policies learned from the Reddit TL;DR summarization task on an out-of-domain dataset: the test partition of the CNN/DailyMail news articles \citep{nallapati-etal-2016-abstractive}, utilizing the optimal sampling temperatures derived from TL;DR (0 and 0.25). Table~\ref{tab:ood} displays the comparative outcomes. We calculated the win rates for GPT-4 versus the gold-standard dataset summaries, applying the same GPT-4 (C) prompt used for Reddit TL;DR, with the term ``forum post'' substituted by ``news article''. On this shifted distribution, DPO consistently exceeds the performance of PPO by a notable gap. These results offer preliminary support that DPO can generalize as effectively as PPO, despite not leveraging the extra unlabeled Reddit TL;DR prompts that PPO uses.

\subsection{Validating GPT-4 judgments with human judgments}
\label{sec:human-judgments}
We undertook a human evaluation to assess the dependability of GPT-4's comparative judgments, employing outputs from the TL;DR summarization study alongside two GPT-4 prompting variants. The \textbf{GPT-4 (S)} (simple) prompt requests a choice of the summary that better covers the key information in a post. By contrast, the \textbf{GPT-4 (C)} (concise) prompt also solicits a preference for conciseness; we examined this variation as we noticed GPT-4 favors longer and more repetitive summaries than human evaluators when following the \textbf{GPT-4 (S)} prompt. Full prompt texts are given in Appendix~\ref{app:prompts}. Our evaluation covered three system comparisons—selecting the best-performing (DPO, temp. 0.25), the worst-performing (PPO, temp. 1.0), and an intermediate (SFT, temp. 0.25) approach—to capture a range of summary qualities; each was benchmarked against greedily-sampled PPO (the best PPO temperature setting). Our analysis shows that, for both prompts, GPT-4's rankings align with human ratings at rates comparable to inter-annotator agreement among human judges, indicating GPT-4 is a credible surrogate for human judgment (given the small number of raters, multiple human assessments were gathered only for DPO and PPO-1 comparisons). In general, the \textbf{GPT-4 (C)} prompt yields win rates that more closely reflect human preferences, motivating its use in our primary experiments in Section~\ref{sec:dpo-real-datasets}. More information about the human evaluation, including the web interface and participant roster, can be found in Appendix~\ref{app:human-study}.

\section{Discussion}
Preference-based learning offers a robust and scalable approach for developing language models that are both effective and aligned with user intent. In this work, we have proposed DPO, a straightforward training framework that leverages preferences to train language models, eliminating the need for reinforcement learning. Instead of recasting the preference learning challenge into a traditional RL framework solely to utilize pre-existing RL techniques, DPO establishes a correspondence between language model policies and reward representations. This connection allows for the direct optimization of language models to meet human preferences via a simple cross-entropy objective, entirely circumventing reinforcement learning without sacrificing expressivity. Remarkably, DPO achieves comparable or superior results to established RLHF approaches—including those that employ PPO—without extensive hyperparameter tuning. As such, DPO offers a more accessible pathway for preference-driven language model training.

\textbf{Limitations \& Future Work.} Several avenues for further research are evident from our findings. One important question concerns the ability of the DPO-trained policy to generalize to distributions distinct from those seen during training, especially when compared to methods that rely on explicit reward functions. Preliminary evidence indicates that DPO’s generalization may match that of PPO-based methods, but broader evaluation is necessary. Another topic is the efficacy of using self-generated labels from DPO in leveraging unlabeled prompts. Moreover, understanding how phenomena like reward over-optimization emerge within the DPO paradigm warrants further scrutiny—such as whether the slight drop in performance observed in Figure~\ref{fig:dialogue-main}-right is an example. While our experiments are conducted on models up to 6B parameters, scaling DPO to models significantly larger and closer to the current state-of-the-art represents a promising direction. Regarding assessments, we observe that GPT-4’s computed win rates are influenced by the prompt itself, suggesting future exploration into methods for reliably obtaining high-fidelity judgments from automated evaluators. Beyond language modeling, DPO may also serve in training generative models across other data modalities, opening up numerous application possibilities.